% ===== §18c: the co-structure of the limit =====
\section{The Co-Structure of the Limit}
\label{sec:costructure}

\noindent
This section gives the limit a formal counterpart. The aim is to show that two figures the paper has described in prose, a lack that carries meaning rather than marking a mere absence, and a core knowable only through its effects on what surrounds it, have precise expressions, the first in the language of coalgebra and the second in the theory of reduced dynamics for open systems. The formal claims are offered as tools for later work and as heuristics that make the prose precise, not as laws to which relational being is asserted to conform. The section states the algebra-coalgebra duality, reifies the lack as a co-structure, gives the semantics of observation and indistinguishability, and then treats rigorously the case of the open composite system in which one component is knowable only through the others.

\subsection*{A note on the ambition of this section}

A word on what is and is not claimed. The formalism below is not offered as evidence that relational being obeys a transcendental law, and the paper does not hold, here or anywhere, that intimacy or drive or value is governed by an equation. The coalgebraic language and the reduction formula are used in the opposite spirit, as heuristics that give a precise home to figures the paper has already established on other grounds, and as a possible basis for later formal work. Where an equation appears, it states a relation that holds in a well-defined mathematical setting and is imported as an analogy whose fit is argued, not as a claim that the setting is the hidden truth of the relational phenomenon. This matches the stance the series takes elsewhere toward optimization and law: a formal structure can illuminate without legislating.

\subsection*{The algebra-coalgebra duality}

The first figure is the meaning-bearing lack, and the language that makes it precise is the duality between algebra and coalgebra. An algebra for a functor $F$ is a map
\[
\alpha : F X \longrightarrow X,
\]
an operation that takes structured collections of elements and combines them into a single element. Algebra is the mathematics of construction and of closure: it says how parts are assembled into a whole, how many are folded into one. A coalgebra for a functor $F$ is the dual, a map
\[
\gamma : X \longrightarrow F X,
\]
an operation that takes a single element and unfolds it into its structured observable behaviour, its one-step outward form together with its successors. Coalgebra is the mathematics of observation and of process: it does not assemble a whole from parts but unfolds a state into what can be seen of it and what it becomes \citep{rutten2000universal,jacobs2016introduction}.

The relevance to the argument is exact. The operation the paper has warned against throughout, the closure of the gap, the folding of the remainder into the signifier, the assembling of the whole into a possessed unity, is an algebraic operation, a map $F X \to X$ that would deliver the many into one and close them. The operation the paper has affirmed, the circling that unfolds without closing, the drive that turns and the desire that slides, the generativity that is never settled, is coalgebraic, a map $X \to F X$ that unfolds a state into its behaviour and its successors without gathering them into a completed whole. The distinction between a language reconciled to its limit and a language that would abolish it is, formally, the distinction between coalgebra and algebra.

\subsection*{The lack reified as a co-structure}

The prose claim that a lack can carry meaning, that attending to any part of a relational structure reveals it as a lackness which is nonetheless meaning-bearing, admits a construction. Take a structure given as a graph of relations, its nodes the terms and its arrows the relations among them. To attend to a particular place is to ask what that place is, and the relational answer is that the place is nothing but the pattern of arrows incident to it. Now perform two operations. First, reify the lack: introduce, as a node of the structure, the very place that was empty, so that the absence becomes a bearer of relations rather than a hole in the diagram. Second, reverse the arrows: pass from the relations that pointed toward a term to the co-relations that unfold from the place. The result is a co-graph in which the reified lack is a node like any other and the arrows run outward from it, so that what was a missing centre becomes a source from which the structure unfolds.

This construction is coalgebraic in form. The reversal of the arrows is the passage from $F X \to X$ to $X \to F X$, from relations that assemble toward a term to co-relations that unfold from a place; and the reification of the lack recognizes the place of absence as a state of the coalgebra, one that unfolds its own outward behaviour. The lack becomes an entity whose whole content is the pattern of what unfolds from it, which is the coalgebraic mode of being, a state known by its observable unfolding and by nothing else.

\subsection*{The semantics of observation and of indistinguishability}

Two notions give the account its semantics. The first is the final coalgebra. Among all coalgebras for a functor there is, under suitable conditions, a final one, into which every other maps uniquely, and it has the character of the space of all possible behaviours, never itself an assembled object but the completed horizon of what can be observed \citep{rutten2000universal}. The final coalgebra is the formal figure of a generativity that is never closed and yet is not formless: the whole of the behaviour without being a whole that could be possessed. It is the coalgebraic counterpart of what the paper has called the open centre, given not as an object but as the totality of what unfolds around a place.

The second is bisimulation, the coalgebraic notion of sameness. Two states are bisimilar when no observation distinguishes them, when everything that can be unfolded from the one can be matched by an unfolding from the other \citep{jacobs2016introduction}. What is decisive is that bisimulation defines identity through observable behaviour alone, without appeal to any internal content the states might be supposed to have. This is the formal correlate of a recurring claim of the paper, that the core is knowable only through its effects. In the coalgebraic setting there is no internal content behind the behaviour to which one might appeal; a state is the pattern of its observable unfolding, and to know it is to observe that pattern, never to seize a content beneath it.

\subsection*{A family of parallels, offered as analogy}

Before the rigorous case, a word on a family of concepts that share the co-structural character loosely rather than by any claimed isomorphism. The signifier of the lack in the Other, the point at which the symbolic order fails to close upon itself, anchors the order not by being a further element of it but by marking the place where an element is missing. One is tempted to place beside it the efficient point of an economy, defined by the trade-offs around it, or a distinguished point of a dynamical flow, or the terms rendered as emptiness and void in the contemplative traditions. These parallels are suggestive, and the paper offers them as analogy and not as identity. They are not shown here to be instances of one mathematical structure, and the section does not rest any weight on their unity. What follows, by contrast, is a case where the correspondence is exact.

\subsection*{The open composite system, treated rigorously}

The rigorous case is the open composite system, and it makes precise the paper's central epistemic claim, that a core is known only through its effects on what surrounds it. The claim is not a metaphor. It is a theorem of the theory of reduced dynamics, and the following states it carefully.

Let the state space of a composite system be the product
\[
M = M_A \times M_B \times M_C ,
\]
with flow $\Phi^t$ generated by a coupled vector field $f$, where coupling means that $f$ does not separate across the three factors. Observables are functions $g \in L^2(M,\mu)$, and they evolve by the Koopman operator
\[
(U^t g)(x) = g(\Phi^t x),
\]
which is linear on observables even when $f$ is nonlinear, this linearity being the point of the Koopman formulation \citep{koopman1931hamiltonian}. Suppose now that only the components $A$ and $C$ are accessible, so that the observables one can actually read form the closed subspace $\mathcal{O}_{AC} \subset L^2(M,\mu)$ of functions depending on $(x_A, x_C)$ alone. Let $\mathcal{P}$ be the orthogonal projection onto $\mathcal{O}_{AC}$, equivalently the conditional expectation $\mathbb{E}[\,\cdot \mid x_A, x_C\,]$ that averages out $B$, and let $\mathcal{Q} = I - \mathcal{P}$. Writing $L$ for the generator of $U^t$, the evolution of the accessible part of any observable satisfies the Mori--Zwanzig generalized Langevin equation
\[
\frac{d}{dt}\, \mathcal{P} U^t g
= \underbrace{\mathcal{P} U^t \mathcal{P} L\, g}_{\text{Markov}}
+ \underbrace{\int_0^t \mathcal{P} U^{t-s} \mathcal{P} L\, e^{s \mathcal{Q} L} \mathcal{Q} L\, g \; ds}_{\text{memory}}
+ \underbrace{\mathcal{P} U^t\, e^{t \mathcal{Q} L} \mathcal{Q} L\, g}_{\text{fluctuation}} ,
\]
an identity that holds exactly, with the memory kernel
\[
K(s) = \mathcal{P} L\, e^{s \mathcal{Q} L} \mathcal{Q} L
\]
carrying the influence of the inaccessible dynamics \citep{zwanzig1961memory,mori1965transport,chorin2000optimal}.

The content of the formula, for the present purpose, is precise. The component $B$ never appears as a direct observation. There is no operator on $B$ in the equation for the accessible dynamics; $B$ has been projected out by $\mathcal{P}$. Yet $B$ is not absent from the dynamics. Its entire effect is carried by the memory term and the fluctuation term, that is, by the marks it leaves on the evolution of the accessible observables over $(A,C)$. To know $B$, in this exact sense, is to read its signature in the memory kernel $K(s)$ that governs how $A$ and $C$ evolve, and there is no other access to it, since the coupled and open structure permits no operator on $B$ itself. The naive picture, on which one would observe $B$ by acting on $B$, is not merely impractical but formally unavailable: the accessible algebra is $\mathcal{O}_{AC}$, and $B$ enters only through what $\mathcal{Q}$ carries into the memory of the accessible part.

An exactly parallel statement holds in the quantum setting, where the accessible observables form the subalgebra $\mathcal{B}(\mathcal{H}_A) \otimes \mathbf{1}_B \otimes \mathcal{B}(\mathcal{H}_C)$, the inaccessible component is removed by the partial trace $\rho_{AC} = \operatorname{tr}_B \rho_{ABC}$, and $B$ survives only in its effect on the reduced state $\rho_{AC}$. The two settings, classical reduced dynamics and quantum restriction to a subalgebra, give one result: a component of an open composite system is knowable only through its neighbours, never through an operator on itself.

\begin{claim}[The core is read through its neighbours]
In an open composite system whose accessible observables are restricted to two of three coupled components, the third component never appears as a direct observation and yet determines the accessible dynamics entirely through the memory kernel of the Mori--Zwanzig equation. This is an exact result, not an analogy, and it is the formal image of the paper's epistemic claim: the unsymbolizable core is the projected-out component, known only through the marks it leaves on the observable neighbours, the symptom, the signifier, and the texture of a shared world, and never through an operator directed at it.
\end{claim}

\subsection*{The reach of the formalism}

Two things are established and one is disclaimed. Established: the meaning-bearing lack has a coalgebraic form, a state whose being is the pattern of what unfolds from it; and the claim that a core is known only through its neighbours is exactly true in the theory of reduced dynamics, where the projected-out component lives on in the memory kernel and nowhere else. Disclaimed: the paper does not assert that relational being is a dynamical system obeying this equation, nor that the family of parallels gathered above are instances of one structure. The formalism is a lamp, not a law. It shows that two of the paper's central figures, the generative lack and the core read through its effects, are not loose metaphors but have precise formal counterparts, and it marks a direction for later work in which the relational ontology of the series might be given a systematically coalgebraic and reduction-theoretic formulation, with closure on the algebraic side, generativity and observation on the coalgebraic side, and the inaccessible core figured throughout as what survives only in the memory of its neighbours.
